В настоящей работе проведён сравнительный анализ систем вёрстки документов {\LaTeX} и Typst. По результатам анализа можно сделать следующие выводы.

\begin{enumerate}
    \item \textbf{Синтаксис.} Typst предлагает значительно более читаемый и лаконичный синтаксис. Документы на Typst в среднем на 30–40\,\% короче эквивалентных документов на {\LaTeX}, что снижает когнитивную нагрузку на автора.

    \item \textbf{Производительность.} Инкрементальный компилятор Typst обеспечивает кратно более высокую скорость: коэффициент ускорения $S$ (формула~(\ref{eq:speedup})) составляет от 7 для коротких документов до более чем 30 для объёмных, что критически важно при итеративной работе.

    \item \textbf{Диагностика ошибок.} Typst выдаёт точные и понятные сообщения об ошибках, тогда как {\LaTeX} нередко генерирует многостраничный вывод с малоинформативными сообщениями.

    \item \textbf{Экосистема.} {\LaTeX} по-прежнему существенно превосходит Typst по зрелости и количеству доступных пакетов \cite{mittelbach2004latex}. Репозиторий CTAN насчитывает тысячи пакетов, тогда как экосистема Typst только формируется.

    \item \textbf{Совместимость.} Многие научные журналы и конференции требуют материалы именно в формате {\LaTeX}. До момента появления официальных шаблонов Typst для целевых изданий переход затруднён.
\end{enumerate}

\textbf{Рекомендация.} Для новых проектов целесообразно рассмотреть Typst как основную систему вёрстки: выигрыш в производительности и удобстве работы перевешивает временную ограниченность экосистемы. Для публикаций в журналах, требующих {\LaTeX}, придется придерживаться сложившейся практики.

Таким образом, Typst представляет собой перспективную альтернативу {\LaTeX} и, по моему мнению, займёт значимое место в инструментарии академических и технических авторов в ближайшие годы.

Чтобы убедиться в этом на практике, достаточно взглянуть на исходный код этого документа: сравните классический синтаксис LaTeX с чистым кодом Typst, выдающими идентичный результат.

Все же, если не мы будем развивать экосистему Typst, приближая будущее с по"=настоящему быстрым и удобным инструментом, то кто сделает это за нас?
